\documentclass{article}
\usepackage{amsmath}
\usepackage{listings}

\begin{document}

\section{Formal Grammars}

A formal grammar consists of variables, terminal symbols, and production rules.
We define a context-free grammar as a tuple with a set $V$ of variables and a set $\Sigma$ of terminals.
Each production rule replaces a variable with a string of symbols.

The derivation of a sentence yields valid strings:
\begin{equation}
    S \implies^* w.
\end{equation}

The yield probability of deriving a sentence is well defined:
\begin{align*}
    S \implies \alpha \implies^* w \quad \text{providng that } w \in L(G).
\end{align*}

In our parsing pipeline, we represent productions using the \verb|A -> B C| syntax or \lstinline{A -> 'token'} rules.
A sample grammar rule file is specified in the following format:
\begin{lstlisting}[language=caml]
S -> NP VP
NP -> Det N
VP -> V NP
\end{lstlisting}

The parsing algorithm analyzes input tokens in $O(n^3)$ time complexity.
The parser calculates the most probable derivation quiickly when the grammar is in Chomsky normal form.

\end{document}
